\documentclass{ctexart}
\usepackage{amsmath,amssymb}
\begin{document}
\section*{固定线路与需求响应交通的两阶段随机规划}
本文件完整定义旧02实际选定的给定场景概率SP模型。
论文标题含DRO，但旧02的目标是$\sum_w\rho_w OBJ_w$；
该加权目标定义给定需求情景上的随机规划。按原文第9--12页修正首站
基准时刻、编组整数域及本场景客流求和。

\subsection*{输入、第一阶段时刻与停站}
站点$U=\{0,\ldots,S-1\}$，列车$I$按输入计划顺序编号，
时间网格$\mathcal T=\{1,\ldots,T_{\max}\}$以分钟计，$\Delta=1$。
乘客组$p\in P_w$有起点$O_p$、终点$D_p>O_p$、到达时刻$AT_p\ge1$、
人数$n_p$。场景概率非负且和为1。模块容量$C$、编组上限$M_f$，
中间站初始模块$N_u^0$与库容$\widehat N_u$有限。
首站基准到达$\bar a_i$与决策$a_{i0}$分开。
\begin{align*}
 a_{i0}&=\bar a_i+\iota_i,&\underline\iota_i\le\iota_i&\le\bar\iota_i,\\
 d_{iu}&=a_{iu}+s_{iu},&
 \underline s_{iu}y_{iu}\le s_{iu}&\le\bar s_{iu}y_{iu},\\
 a_{i,u+1}&=d_{iu}+r_{iu},&
 r_{iu}&=\sum_t\widehat r_{u,u+1,t}h_{iut}
            -r^{acc}_{iu}(1-y_{iu})-r^{dec}_{i,u+1}(1-y_{i,u+1}),\\
 d_{i+1,u}-d_{iu}&\ge h,& y_{i0}=y_{i,S-1}&=1.
\end{align*}
$a,d,s$为非负整数分钟，$y$为二元停站变量。
$e_{iut}=\mathbf1_{\{t<a_{iu}\}}$、$z_{iut}=\mathbf1_{\{t<d_{iu}\}}$，
$g_{iut}=\mathbf1_{\{t=a_{iu}\}}$、$h_{iut}=\mathbf1_{\{t=d_{iu}\}}$。
故$e,z$随$t$非增，$g_{iu1}=1-e_{iu1}$、
$g_{iut}=e_{iu,t-1}-e_{iut}$，$h$同理；
$a_{iu}=1+\sum_te_{iut}$，$d_{iu}=1+\sum_tz_{iut}$。
当前输入$\underline s=\bar s=$既有停站时长，路段运行时间恒定；
跳站减速/加速节约作为缺项取0。算法选择所有站停靠，完整模型仍允许中间站跳站。
时间域取足以覆盖全部候选及行程结束的显式有限整数区间。

\subsection*{第二阶段服务选择、容量}
FLS编组$x_{iu}^w=\sum_{m=1}^{M_f}m q_{ium}^w$，
$\sum_mq_{ium}^w=1$，$q$二元，因此$x\in\{1,\ldots,M_f\}$。
DRT变量$\pi_{uvtm}^w$二元，$\eta_{uvt}^w=\sum_m\pi_{uvtm}^w\in\{0,1\}$。
DRT行程时间$\widehat r_{uvt}>0$；当前数据为$(v-u)r$。
组$p$的可用服务为停靠其起终点、$d_{i,O_p}\ge AT_p$的FLS，
以及$\eta_{O_p,D_p,t}^w=1,\ t\ge AT_p$的DRT。
必须有可用服务。每组恰好分配给其中发车最早的一班；
沿用旧02严格选择约束，若最早时刻存在并列，该候选不允许。
令二元$\varphi_{pk}^w$表示该选择，则
\[
\sum_{k\in A_p^w}\varphi_{pk}^w=1,\quad
WT_p^w=\min_{k\in A_p^w}(\operatorname{dep}_k-AT_p),\quad
\varphi_{pk}^w=1\Longrightarrow
 \operatorname{dep}_k-AT_p=WT_p^w,
\]
且所有未选可用服务的等待时间严格更大（整数分钟下至少大1）。
$TP_p^w=\sum_k\varphi_{pk}^w\operatorname{ride}_{pk}$，
FLS乘车时间$a_{i,D_p}-d_{i,O_p}$，DRT为$\widehat r_{O_p,D_p,t}$。
不使用服务的候选等待时间取足够大的$M$，对应旧02的辅助变量。
\[
 b^w_{iu}=\sum_{p:O_p=u}n_p\varphi_{pi}^w,\quad
 l^w_{iu}=\sum_{p:D_p=u}n_p\varphi_{pi}^w,\quad
 o^w_{iu}=\sum_{p:O_p\le u<D_p}n_p\varphi_{pi}^w
 \le Cx^w_{iu}.
\]
末站载客量为0。DRT人数
$b^w_{uvt}=\sum_{p:(O_p,D_p)=(u,v)}n_p\varphi^w_{p,uvt}
\le C\sum_m m\pi^w_{uvtm}$。

\subsection*{有限车辆库存}
每个中间站$1\le u\le S-2$，所有事件同一时刻先按净变化合并：
\begin{align*}
 N^w_{ut}&=N^w_{u,t-1}+AV^w_{ut}-DV^w_{ut},&
 0\le N^w_{ut}&\le\widehat N_u,\qquad N^w_{u0}=N_u^0,\\
 AV^w_{ut}&=\sum_i x^w_{i,u-1}g_{iut}
 +\sum_{\substack{v,\tau,m:\\\tau+\widehat r_{vu\tau}=t}}m\pi^w_{vu\tau m},\\
 DV^w_{ut}&=\sum_i x^w_{iu}h_{iut}+\sum_{v,m}m\pi^w_{uvtm}.
\end{align*}
始终使用实际出发$\tau$对应的行程时间，避免旧02的反向时耗索引。
不存在的出发变量为0。按旧02，库存限制仅在中间站；
两端站供应不受上述中间站库容约束。

\subsection*{目标与算法}
\[
\min\sum_w\rho_w\left\{
\zeta_{pa}\varphi_T\sum_{p\in P_w}n_p(\xi_1WT_p^w+\xi_2TP_p^w)
+\zeta_{op}\left[\sum_{i,u,m}\varphi_{u,u+1,m}q^w_{ium}
+\sum_{u,v,t,m}\varphi_{uvm}\pi^w_{uvtm}\right]\right\}.
\]
当前补全运行费$\varphi_{uvm}=(0.55+0.35m)\widehat r_{uv}$，
时间价值0.25/分钟，等待/乘车权重2/1；均按实例保留的数值。
代码选择一个完整全停站时刻表，各场景先生成FLS可行候选，再插入改善DRT；
候选构造方法采用列车全程同一编组。
返回真实加权费用及全部客流/库存证书，只是SP可行上界。
结果记录为SP的可行上界。
每实例实际预算3600秒，全部实例由当前入口重跑。
\end{document}
